\documentclass{article}

% Language setting
% Replace `english' with e.g. `spanish' to change the document language
\usepackage[english]{babel}

% Set page size and margins
% Replace `letterpaper' with `a4paper' for UK/EU standard size
\usepackage[letterpaper,top=2cm,bottom=2cm,left=3cm,right=3cm,marginparwidth=1.75cm]{geometry}

% Useful packages
\usepackage{amsmath,amsfonts,amssymb}
\usepackage[colorlinks=true, allcolors=blue]{hyperref}

% Define theorem like objects
\newtheorem{theorem}{Theorem}

\usepackage{mathtools}
\mathtoolsset{showonlyrefs}

\newcommand{\e}{\varepsilon}
\newcommand{\dx}{\, dx}
\newcommand{\dt}{\, dt}
\newcommand{\ds}{\, ds}
\newcommand{\dy}{\, dy}
\newcommand{\dz}{\, dz}
\newcommand{\dr}{\, dr}
\newcommand{\x}{{ \bf x}}
\newcommand{\w}{{ \bf w}}
\newcommand{\al}{\alpha}
\newcommand{\R}{\mathbb{R}}
\newcommand{\N}{\mathbb{N}}
\newcommand{\E}{\mathbb{E}}
\newcommand{\HH}{\mathcal{H}}
\DeclareMathOperator{\Prob}{\mathbb{P}}
\renewcommand{\P}{\mathbb{P}}
\newcommand{\M}{\mathcal{M}}
\renewcommand{\div}{\text{div}\,}
\newcommand{\sym}{\text{sym}}
\newcommand{\diam}{\textup{diam}\,}
\newcommand{\grad}{\nabla}
\newcommand{\lambdaMin}{\lambda_{\textup{min}}}
\newcommand{\lambdaMax}{\lambda_{\textup{max}}}
\newcommand{\vxavg}{\vx_{\textup{avg}}}
\newcommand{\xavg}{x_{\textup{avg}}}
\newcommand{\xperp}{x_{nc}}
\newcommand{\vxperp}{\vx_{nc}}
\newcommand{\ones}{{\bf 1}}
\newcommand{\myDiag}{\textup{diag}\,}
\newcommand{\ddt}{\frac{d}{dt}}
\newcommand{\cl}{\textup{cl}}
\newcommand{\dtau}{\,d\tau}
\newcommand{\myspan}{\textup{span}}
\newcommand{\Var}{\textup{Var}}
%\newcommand{\C}{\mathfrak{C}}
\newcommand{\C}{\mathcal{C}}
\newcommand{\cons}{\mathfrak{C}}
\newcommand{\vecc}{\text{vecc}}
\newcommand{\diag}{\text{diag}}
\newcommand{\ns}{{n_s}}
\newcommand{\Proj}{\textup{P}}
\newcommand{\Hess}{\nabla^2}
\newcommand{\myrank}{\mbox{rank}}
\newcommand{\ELBO}{\text{ELBO}}
\newcommand{\calL}{\mathcal{L}}
\newcommand{\calF}{\mathcal{F}}
\newcommand{\eps}{\varepsilon}

% KL divergence definition
\DeclarePairedDelimiterX{\KLdivx}[2]{(}{)}{%
  #1\;\delimsize\|\;#2%
}
\newcommand{\KLdiv}{D_{\text{KL}}\KLdivx}
\DeclarePairedDelimiter{\norm}{\lVert}{\rVert}

\title{VAE Notes}
\author{}
\date{}

\begin{document}
\maketitle


\section{Setup}

\noindent \textbf{High level Goal}: We have some dataset $D= (x_i)\subset X$. We assume that $D$ was generated by iid sampling from some distribution $p(x)$. We'd like to model $p(x)$. \\

\noindent \textbf{Motivation}:
Consider the following example. Let $z\in \R^\ell$ be some latent variable, where the dimension $\ell$ is a hyperparameter. Let $p_\theta(x, z)$ be a parametric joint density function over $x$ and $z$. Let $q_\phi(z|x)$ be a conditional density over $z$ conditioned on $x$. We do not assume these are intrinsically coupled or tied to the true density $p$ in any way (yet). They are simply parametric density functions. Solve the optimization problem
\begin{equation} 
\max_{\theta, \phi} \sum_{x\in D} \E_{z\sim q_{\phi}(z|x)}\left[ \log \frac{p_\theta(z, x)}{q_\phi(z|x)} \right].
\end{equation}
See notebook example. This is done with exactly the MNIST data set. There is a bit of art to choosing the form of the parametric densities. But note the continuous properties of the latent space. Note the reconstruction error. This architecture is called a variational autoencoder. These notes derive the variational autoencoder, hopefully clarifying where this loss comes from, how this is architecturally related to an autoencoder, and hopefully giving insights on how they work and when to use them. 

Other motivating examples: consider latent diffusion models. In particular stable diffusion XL. Show example samples. And then show a diagram of the architecture, with a VAE used to produce the latent space.  Review graph-based semi-supervised learning techniques. VAE is used to produce latent space where nice similarity graph can be constructed. Review automatic chemical design paper. Really cool thing about this is that allows us to do discrete design by moving into a well structured continuous latent space. Design in the continuous latent space is a lot easier. (What kinds of design things can you do in the latent space?) 


% As a third example, could consider semi-supervised learning examples where a VAE is used to produce a latent space where a similarity graph can meaningfully be constructed. 

\vspace{1em}
\noindent \textbf{VAEs}
There are a lot of ways we could try to model $p$. VAEs are often motivated from the perspective of using a graphical model. I'm sure there's a good reason for this. But it doesn't make a lot of sense to me. My take on it is the following: We would like to try and approximate the true distribution $p(x)$ with a latent variable model
$$
p_\theta(x, z) = p_\theta(x|z)p_\theta(z)
$$
in the sense that the marginal distribution $p_\theta(x) = \int p_\theta(x,z)\dz \approx p(x)$. Typically, the space $X$ is some Euclidean space and the latent space $Z$ is a lower dimensional Euclidean space. Why use a latent variable model? In some sense, this allows you to implicitly reduce the degrees of freedom of your model. The manifold hypothesis presumes that data is distributed on some low dimensional manifold in a high dimensional ambient space. This allows you to try to pick your latent dimension to match the intrinsic dimension of the distribution. 
(Other reasons for using a latent variable model? Interpetability?)


% As a tangent, it seems worth noting that given our dataset $D$ alone, the only approximation we can come up with is the empirical distribution on $D$. Any other distribution we construct to approximate $p$ leverages our prior assumed knowledge or inductive bias in some way. Some of the miracle of generative modeling techniques to me is that they somehow implicitly regularize in the right way to promote good generalization. 


\vspace{1em}
\noindent Rough outline: 
\begin{itemize}
  \item Try to apply variational inference to solve the problem. 
  \item When we run into computational bottlenecks, plug a neural network in at the right place. This will give us ``amortized'' variational inference, which amortizes the computational cost.
  \item To actually train the neural network, we'll want to be able to differentiate with respect to our parameters. This will be tricky, because of the stochastic elements we've introduced. To compute derivatives, we'll use the ``reparameterization trick''. In the original VAE paper, they suggest that this trick, rather than the VAE itself, could be considered the main contribution. So, it's worth being aware of. 
  \item When you put those together, you get a neural network architecture for $p_\theta(x, z)$ and a loss function, differentiable in $\theta$, that you can use for training. 
  \item I have some computational examples that show how it fits together.

\end{itemize}

\section{A Variational Inference Approach}
% In my experience, (mostly just this paper \cite{blei2017variational},) variational inference (VI) is traditionally used in Bayesian inference to approximate a tough-to-compute posterior $p(z|x)$, where $z$ represents some latent variables in an assumed latent variable model $p(x,z)$. Here, the objective is slightly different. We're just trying to approximate $p(x)$. And we're going to use a latent variable model to do it. I note this because if you try and compare the literature on VI in Baysian inference and VAEs, it gets confusing quickly because they both use latent variable models and VI techniques, but applied in different ways and with different rationale. With that in mind, I'll discuss how they apply VI techniques to come up with the VAE. 

Our goal is to approximate $p(x)$. Let $q$ be some distribution defined over the latent space $Z$. For now, leave $q$ unspecified. After we set up objectives, we'll be able to optimize over $q$, which is the ``variational'' part of variational inference.\footnote{Classically, with CAVI type stuff (section 2.4 of \cite{blei2017variational}), I think they really did solve some variational problems as part of VI. But the more modern stuff I've seen just paramterizes things with a neural network and then optimizes using GD in parameter space. But I'm not super familiar with the literature.} Let $x\in X$ and observe that
\begin{align}
    \log p_\theta(x) & = \E_{z\sim q}\left[ \log p_\theta(x) \right]\\
    & = \E_{z\sim q}\left[\log \frac{p_\theta(x,z)}{p_\theta(z|x)} \right]\\
    & = \E_{z\sim q}\left[ \log \frac{p_\theta(x,z)}{q(z)} \frac{q(z)}{p_\theta(z|x)} \right]\\
    & = \underbrace{\E_{z\sim q}\left[ \log\frac{p_\theta(x,z)}{q(z)}\right]}_{:= \mbox{ELBO}(\theta, q)} + \underbrace{\E_{z\sim q}\left[\log \frac{q(z)}{p_\theta(z|x)} \right]}_{=\KLdiv{q}{p_\theta(z|x)}} \label{eq:ELBO-pls-KL}
\end{align}
Because $D_{KL} \geq 0$, the ELBO is a lower bound on $p_\theta(x)$
$$
p_\theta(x) \geq ELBO(\theta, q).
$$
In the ML literature, the ELBO is also commonly referred to as the ``variational lower bound'' and sometimes ``variational free energy''. This can also be derived using Jensen's inequality. 

\vspace{1em}\noindent 
A few thoughts on interpreting the ELBO: 
\begin{enumerate}
\item  Since $p_\theta(x)$ is constant ($x$ is fixed as far as any optimization is concerned), maximizing $\ELBO(\theta, q)$ is exactly the same as minimizing $\KLdiv{q}{p_\theta(z|x)}$. Hence, we maximize the ELBO (as will be our objective) this can be thought of as equivalently driving these two distributions together.

\item Following from the previous point, for a fixed $x$ and $\theta$, the bound becomes tight if $q = p_\theta(z|x)$.
\end{enumerate}

3. For an alternate perspective on what it means to maximize the ELBO, massage things a bit to get
\begin{align}
\ELBO(\theta, q) & = E_q\left[ \log \frac{p_\theta(x,z)}{q(z)} \right]\\
& = \E_q\left[ \log p_\theta(x|z) \right] + \E_q[\log p_\theta(z)]  + E_q[\log q(z)]\\
& = \E_q[\log p_\theta(x|z)] + E_q[\log \frac{p_\theta(z)}{q(z)}]\\
& = \underbrace{\under\E_q[\log p_\theta(x|z)]}_{(1)} - \underbrace{\KLdiv{q(z)}{p_\theta(z)}}_{(2)}
\end{align}
To build intuition, suppose that $p_\theta(x|z)$ has the form
$$
p_\theta(x|z) = N(x; G_\theta(z), \eta I),
$$
meaning a normal distribution over $x$ with mean $G_\theta(z)$, and covariance $\eta I$. Here, we'll let $G_\theta$ be a network and $\eta, \theta$ be parameters. 

We look at (1) above, first. Inside the expectation we have
\begin{align}
    \log p_\theta(x|z) & = \log N(x; G_\theta(z), \eta I)\\
    & = \log \left[ (2\pi \eta)^{-d/2} \exp(-\frac{1}{2}\|x - G_\theta(z)\|^2) \right]\\
    & = \frac{1}{2\eta} \|x - G_\theta(z)\|^2 + \text{const}.
\end{align}
This is like mean squared reconstruction error. For fixed $\theta$, this encourages the mass of $q$ to be placed on $\arg\min_z \|x - G_\theta(z)\|$. 

Understanding (2). Typically, $p_\theta(z)$ is fixed to be $N(z; 0, I)$, with no dependence on $\theta$. This term encourages $q$ to be close to the "prior" $p(z)$. The interpretation for fixed $q$ and varying $\theta$ is less clear.

Comments: 
\begin{itemize}
    \item I find this particularly confusing. Because the ELBO loss can clearly be decomposed into the two components above. The second component clearly drives $q(z)$ towards the prior $p_\theta(z)$, which I'm going to think of as a unit normal. \emph{But}, \eqref{eq:ELBO-pls-KL} clearly shows that maximizing the ELBO is equivalent to minimizing the KL divergence between $q(z)$ and the posterior $p_\theta(z|x)$. So, ignoring the decomposition of the ELBO, at maxima, the LK between $q$ and the posterior are minimized! But this conflicts with the apparent goal in the loss that drives $q(z)$ towards the fixed normal prior. 
    \item The loss here decomposes into a reconstruction loss term and a KL divergence term. The reconstruction loss is MSE when the output density is a parametric normal. Just a note that this seems like a terrible reconstruction loss metric. Something like a Wasserstein distance seems like it would be better. Is there a way to massage this formulation so you get a Wasserstein distance in the reconstruction loss?
    \item Continuing with the previous comment, what is the interpretation if you use a multivariable Bernoulli $p_\theta(x|z)$? Does this seem as obviously terrible a reconstruction loss as the MSE? Does it seem better in some way? 
    \item Continuing again, when you look at the decomposition of the ELBO into a MSE reconstruction loss and a KL divergence term which clearly is inconsistent with the original VAE objective (it drives $q$ towards the prior instead of the posterior??), the ELBO seems like a kind of crappy surrogate loss. 
\end{itemize}

\section{Amortized Variational Inference}
For the remainder of this discussion, we will assume that $p_\theta(z) = p(z)$, which is some predetermined density (usually a multivariate standard normal). So, when optimizing $p_\theta(x,z)$, we just optimize the parameters of $p_\theta(x|z)$. 

Given a dataset $D = (x_i)$, we want to solve 
$$
\max_\theta \sum_i \log p_\theta(x_i),
$$
i.e., maximize the likelihood. This is hard, so we optimize the ELBO. Let
$$
\calL_{\theta, q}(x_i) := \E_{z\sim q}\left[ \log\frac{p_\theta(x,z)}{q(z)}\right]
$$
denote the single-data-point ELBO. For a given $x\in D$, let $q_x$ denote a distinct choice of $q$ associated with this $x$. We wish to optimize
$$
\max_{\theta, (q_x)_{x\in D}} \sum_{x\in D} \calL_{\theta, q_x}(x).
$$
Consider the following algorithm.\\

\noindent Repeat forever:
\begin{itemize}
    \item For each $x \in D$: 
    \begin{enumerate}
        \item Fit $q$ to approximate $p_\theta(z|x)$
        \item Update $\theta$ using gradient ascent on ELBO
    \end{enumerate}
\end{itemize}
Observations:
\begin{itemize}
    \item Recall, ELBO is tight for fixed $x$ when $q(z) = p_\theta(z|x)$. Hence, we're maximizing ELBO w.r.t. $q$ in step 1. 
    \item We can fit $q$ in step 1 using some parametric model and gradient ascent, per usual in neural network land.
\end{itemize}

\vspace{1em}
\noindent Problem with this: It's expensive! We have to fit a separate $q$ for each $x$. 

\vspace{1em}
\noindent 
VAE solution: Use a neural network to parameterize a distribution $q_\phi(z|x)$, where this is a distribution on $Z$ space, parameterized by the input $x$ (and the tuning parameters $\phi$). For example, $q_\phi(z|x)$ may be a normal distribution over $Z$ whose mean and covariance are determined by the outputs of a neural network with input $x$ and parameters $\phi$. 

Now, the optimization problem becomes
$$
\max_{\theta, \phi} \underbrace{\sum_{x\in D} \calL_{\theta, \phi}(x_i)}_{:= \mathcal{F}(\theta, \phi)},
$$
where $\calL_{\theta, \phi}(x) := \ELBO(\theta, q_\theta(x); x)$ (this means the single data point ELBO at point $x$). To be explicit, our final optimization problem is
\begin{equation} \label{eq:VAE-objective}
\max_{\theta, \phi} \sum_{x\in D} \E_{z\sim q_{\phi}(z|x)}\left[ \log \frac{p_\theta(z, x)}{q_\phi(z|x)} \right].
\end{equation}

Note: In the VAE paper, they call the general approach variational Bayes, or something, and the VAE is a special case. I don't remember what the distinction is and I may be conflating the general case and the VAE to some degree. 

As a matter of nomenclature, they refer to the technique of using a single neural network to approximate the optimal $q_x$ for all $x$ as \emph{amortized} variational inference. 

\section{Comments}
\begin{itemize}
    \item How do you choose the forms for $p_\theta(x|z)$, $p_\theta(z)$ and $q_\phi(z|x)$?
    \item Some options: $p_\theta(x|z)$ and $q_\phi(z|x)$ both normal distributions with mean and covariance determined by a neural network. 
    \item Consider an application modeling black and white images (e.g., MNIST). $p_\theta(x|z)$ could be a multivariable bernoulii (each pixel is independently represented by a bernoulli, the mean of which is parameterized by a neural network). 
    \item I don't remember if I've already said this, but forcing $p_\theta(x|z)$ to be a neural network with a density function tacked on the end is kind of weird and potentially counterproductive in certain ways. Supposedly, we are using a latent variable model because we think the intrinsic dimensionality of the data manifold is smaller than the ambient space. I have assumed that you'd want the dimension of your latent space to be identical to the manifold dimension. In that case, a continuous mapping from $R^\ell$ to $\R^d$, where $\ell \ll d$ is the manifold dimension, will produce a density function with support on a $\ell$ dimensional subset of $\R^d$. (Whether or not it's a true manifold will depend on differentiability assumptions of the map. Can it be kinky?) By tacking a density to the end of a neural network you're forcing the support of $p_\theta(x|z)$ to have a nonsingular support in the sense that the support is not a set with zero measure.) Which, I imagine, is good for learning (getting gradient information). But not great for truly modeling the data density. Similar ideas are reflected in the Wasserstein GANs paper. Does learning a density with measure-zero support require something like a Wasserstein distance? (TODO: Read Wasserstein autoencoder paper.) 
    \item On the topic of Wasserstein distances, there are two places an earth mover type distance might be useful here. One is in learning a density with singular support. The other is for learning a meaningful similarity distance between images (assuming that's the application). An L2 (MSE) distance between images is an awful similarity metric for the same reason it's a bad distance metric for probability densities. It doesn't handle translation of images.) Reading the Wasserstein VAE paper, I'd like to see if/how they account for each of these. If they don't account for Wasserstein distance in the image distance, is there any way to do that, and would it actually be useful?
\end{itemize}

\section{Optimizing the ELBO and the Reparameterization Trick}
We now consider how to optimize the objective in \eqref{eq:VAE-objective} using gradient ascent (GA). Note that
\begin{align}
\grad_\theta \calF(\theta, \phi) & = \sum_{x\in D} \grad_\theta \E_{z\sim q_\phi}\left[ \log p_\theta(x|z)p(z)\right] - \underbrace{\grad_\theta\E_{z\sim q}\left[ \log q_\phi\theta(z|x)\right]}_{=0}\\ 
& =  \E_{z\sim q_\phi}\left[ \sum_{x\in D}\grad_\theta \log p_\theta(x|z)p(z)\right],
\end{align}
where we assume the necessary assumptions hold to bring the the derivative inside the expectation.
This says that $\sum_{x\in D} \grad_\theta \log p_\theta(x|z)p(z)$ is an unbiased estimator of $\grad_\theta \calF$ when $z\sim q_\theta(z|x)$ for each $x$. Hence, we can optimize in this parameter using SGD (or SGA, I guess). 

Computing the gradient with respect to $\phi$ is harder. Consider computing a gradient of $\calL(\theta, \phi)$ Attempting the same reasoning, note that
\begin{align}
    \grad_\phi \calF(\theta, \phi) = \sum_{x\in D} \grad_\theta \E_{z\sim q_\phi }\left[ \frac{p_\theta(x|z)p(z)}{q_\phi(z|x)} \right].
\end{align}
We can't bring the gradient inside the expectation here. 

\subsection{General Framework for Reparameterization Trick}
To address this, consider the following more general problem. (This comes up in other settings, so it's worth considering it in a general framework.) Suppose we are interested in computing the following gradient
\begin{equation} \label{eq:reparam-grad-goal}
\grad_\phi \underbrace{\E_{z\sim q_\phi(z)} f(z, \phi)}_{:= L(\phi)}.
\end{equation}
Here, $f(\cdot, \phi):\R^n \to \R^n$ is some function (assumptions?) and $\phi$ is the parameter we want to optimize over. We leave out functional dependence on the data point $x$ here, but I don't think that matters. There are typically two ways of addressing this in the ML community/literature. The first is what's known as the likelihood ratio gradient. It's computed as follows
\begin{align}
    \grad_\phi \E_{z\sim q} f(z,\phi) & = \grad_\phi \int q_\phi(z) f(z,\phi)\dz \\
    & = \int \grad_\phi\left( q_\phi(z) f(z,\phi) \right) \dz\\
    & = \int \frac{\grad_\phi\left( q_\phi(z) f(z,\phi) \right)}{q_\phi(z) f(z,\phi)}q_\phi(z) f(z,\phi) \dz\\
    & = \E_{z\sim q_\phi}\left[ \grad_\phi \log\left( q_\phi(z) f(z,\phi) \right) f(z,\phi) \right].
\end{align}
If we sample $z\sim q_\phi$, then $\grad_\phi \log\left( q_\phi(z) f(z,\phi) \right) f(z,\phi)$ is an unbiased estimator of the gradient of interest. Note this is basically the same as the REINFORCE trick from RL. 

\vspace{1em}
\noindent \textbf{Problem}: This has high variance. I have a demo for this. \\


The second method to address this is with the reparameterization trick. For this, you reparameterize your expectation so it's no longer with respect to a distribution which depends on $\phi$. 

To do this, we'll let $p(\eps)$ (represent some distribution over $\R^n$, where we recall $n$ is the dimension of the space where $z$ lives (see definition of $f$ at the beginning of the section).\footnote{$p(\eps)$ conflicts with notation for $p$ as target distribution for VAE.}
Suppose we have a function $g(\eps, \theta)$ so that if 
$$
w = g(\eps, \theta), \quad \eps \sim p(\eps),
$$
then $w$ has density $q_\phi$. Here, $q_\phi$ is the density of $z$ in the expectation in \eqref{eq:reparam-grad-goal}. Note that this density is defined over $\R^n$, which is where $w$ lives. If we can construct such a function $g$ and distribution $p(\epsilon)$, then we have
\begin{align} \label{eq:reparam1}
    \grad_\phi \E_{z\sim q_\phi} f(z, \phi) & = \grad_\phi \E_{\eps \sim p(\eps)} f(g(\eps, \phi), \phi)\\
    & = \E_{\eps \sim p(\eps)} \grad_\phi f(g(\eps, \phi), \phi).
\end{align}
Which gives a way of constructing an unbiased estimator of the gradient of interest. 
Constructing such a distribution $p(\eps)$ and function $g$ is a matter of applying the change of variables formula. 

Note, in \eqref{eq:reparam1}, we were hand wavey about the assumptions required to bring the gradient inside the expectation. I don't remember the exact conditions, but I think they're pretty mild. 


\subsection{Change of Variables}
The change of variables theorem is stated as follows. 
\begin{theorem}[Change of Variables]
    Let $X$ be a random variable with distribution $f_X(x)$. Let $g$ be a differentiable and invertible function with $g^{-1}:=h$ also differentiable (i.e., $g$ a diffeomorphism). Let $Y= g(X)$ and let $f_Y$ denote the distribution of $Y$. Then
    $$
f_Y(y) = f_X(h(y)) |\det D h(y)|.
    $$
\end{theorem}
This is surprisingly difficult to prove. 

\subsection{Example: Reparameterization Trick for Simple VAE}
There is some art to picking a good $g$ and $p(\epsilon)$. A common choice in (simple) VAEs is to represent the encoder distribution $q_\phi(z|x)$ as a normal distribution with mean and covariance given by the output of a neural network. Formally, 
$$
q_\phi(z|x) = N(z; \mu, \diag(\sigma^2),
$$
where $(\mu, \log \sigma) = \mbox{EncoderNet}_\phi(x)$ and $\mu, \sigma \in \R^n$. 
We have
\begin{align}
q_\phi(z|x) & = \prod_{i=1}^n q_\phi(z_i|x)\\
& = \prod_{i} \frac{1}{\sqrt{2\pi}\sigma_i} \exp(\sigma^{-2}(z_i - \mu)^2).
\label{eq:reparam2}
\end{align}
Consider the reparameterization 
\begin{align}
    \eps & \sim  N(0, I)\\
    g(\eps, \mu, \sigma) & = \mu + \sigma \odot \eps,
\end{align}
% By change of variables, we have $f_w(w) = f_\eps(g^{-1}(w)) \left| \det D_w g^{-1}(w)\right|$. Looking at individual terms, $h = g^{-1}(w) = (w-\mu)./\sigma$ and
% $$
% \frac{\partial h_i}{\partial w_j} = 
% \begin{cases}
%     \sigma_i & \quad i = j\\
%     0 & \quad \mbox{else}
% \end{cases}
% $$
where $\odot$ denotes elementwise multiplication. 
Letting $h = g^{-1}$, we have $D_w h(w) = \diag((1/\sigma_i)_i)$ and $\det D_w h = \prod_i \sigma_i^{-1}$. By construction, $f_\eps(\eps) = \prod_{i} N(\eps; 0, 1)$ so
$$
f_\eps(g^{-1}(w)) = \prod_i \frac{1}{\sqrt{2\pi}}\exp{\left(\frac{1}{2}\left(\frac{w_i - \mu}{\sigma_i}\right)^2\right)}
$$
Applying the change of variables formula, we have
$$
f_w(w) = \left(\frac{1}{(2\pi)^{n/2}} \prod_i \exp(\frac{1}{2}\left(\frac{w_i - \mu}{\sigma_i}\right)^2)\right)\prod_i \sigma_i^{-1},
$$
which is the desired density $q_\phi$ in \eqref{eq:VAE-objective}. 

(This next part assumes you're evaluating the loss by decomposing as the numerator/denominator in \eqref{eq:VAE-objective}. If, instead, you decompose into something like a reconstruction term and KL term, then you just sample according to the reparameterization trick, and then pass the sampled latent into the computational graph and use whatever loss terms arise in your decomposition.) In a VAE, we want to apply the reparameterization trick to compute
\begin{align}
    \grad_\phi \tilde L(\phi) & := \E_{z\sim q_\phi(z|x)} \log q_\phi(z|x)\\
    & = \grad_\phi \E_{\eps \sim p(\eps)} \log q_\phi(g(\eps, \mu_\phi, \sigma_\phi)|x)
\end{align}
If we have this, then computing a gradient in \eqref{eq:VAE-objective} is easy. 
Bringing the gradient inside the expectation we have that
$$
\grad_\phi \log q_\phi(g(\eps, \mu_\phi, \sigma_\phi)|x), \quad \eps \sim p(\eps)
$$
is an unbiased estimator of $\nabla_\phi \tilde L$. This is something we can actually evaluate on a computer. The form of $q_\phi$ is given in \eqref{eq:reparam2} and $\mu_\phi$, $\sigma_\phi$ come from a neural network. Hence, we can use the chain rule/autodiff to compute it.  


\section{Example}
See MNIST examples. 


\vspace{1em}
\noindent 
Latent space interpolation: 
\begin{itemize}
    \item \cite{roberts2018hierarchical} https://magenta.tensorflow.org/music-vae
    \item \cite{berthelot2018understanding} 
\end{itemize} 

\section{Questions}
How is latent space structure different for VAE vs autoencoder? 

\section{VAEs in Modern ML}

\noindent \textbf{Image generation:} 
VAEs set the state of the art in generative modeling for images in \cite{razavi2019generating} in 2019. It's no longer state of the art, but the citation count is still growing. 

\vspace{1em}
\noindent \textbf{Image Compression}:
VAEs are closely related to the compression techniques in \cite{balle2017end}.

\vspace{1em}
\noindent \textbf{Interpretable ML}:
Learning Interpretable Latent Space Representations. See $\beta$-VAE \cite{higgins2016beta}. Encourages representation where varying any single latent component has meaningful interpretability. 

\vspace{1em}
Molecule/chemical design: \cite{gomez2018automatic}. (This is similar to the active learning for SAR paper. Train a VAE latent space. Use that for chemical design. e.g., interpolate between chemicals you want. Or design chemicals based on certain properties.) 

\vspace{1em}
Semi-supervised learning: \cite{miller2022graph}. (Setting: Only some of the data has labels, and acquiring new labels is expensive. Train a VAE, unsupervised, on data with no labels. Then use Laplace learning/graph-based methods for semi-supervised learning in the VAE latent space.) 

%https://pubs.acs.org/doi/pdf/10.1021/acscentsci.7b00572

\vspace{1em}
\noindent \textbf{Graphs}: Latent representnation of graphs/link prediction \cite{kipf2016variational}.

\vspace{1em}
\noindent \textbf{Latent diffusion models}

\vspace{1em}
\noindent \textbf{Wasserstein autoencoder} 


\bibliographystyle{alpha}
\bibliography{sample}

\end{document}